\documentclass[11pt,a4paper]{article}

% ============================================================
%  Encodage et langue
% ============================================================
\usepackage[utf8]{inputenc}
\usepackage[T1]{fontenc}
\usepackage[french]{babel}
\usepackage{lmodern}

% ============================================================
%  Mise en page
% ============================================================
\usepackage[a4paper,margin=2.5cm]{geometry}
\usepackage{setspace}
\onehalfspacing
\usepackage{parskip}
\usepackage{microtype}

% ============================================================
%  Math, tableaux, listings, couleurs, liens
% ============================================================
\usepackage{amsmath,amssymb,amsthm}
\usepackage{array}
\usepackage{booktabs}
\usepackage{longtable}
\usepackage{tabularx}
\usepackage{multirow}
\usepackage{graphicx}
\usepackage{xcolor}
\usepackage{listings}
\usepackage{enumitem}
\usepackage{float}
\usepackage{xspace}
\usepackage[hidelinks,colorlinks=true,linkcolor=blue!60!black,
            citecolor=blue!60!black,urlcolor=blue!60!black]{hyperref}
\usepackage{fancyhdr}
\usepackage{titlesec}

% ============================================================
%  Couleurs et style des listings
% ============================================================
\definecolor{codebg}{RGB}{248,248,248}
\definecolor{codekw}{RGB}{0,90,160}
\definecolor{codecm}{RGB}{0,128,0}
\definecolor{codestr}{RGB}{160,40,40}
\definecolor{codeline}{RGB}{120,120,120}
\definecolor{tableheader}{RGB}{230,236,245}

\lstset{
  basicstyle=\ttfamily\small,
  backgroundcolor=\color{codebg},
  keywordstyle=\color{codekw}\bfseries,
  commentstyle=\color{codecm}\itshape,
  stringstyle=\color{codestr},
  numbers=left,
  numberstyle=\tiny\color{codeline},
  numbersep=8pt,
  frame=single,
  rulecolor=\color{codeline!50},
  breaklines=true,
  breakatwhitespace=true,
  tabsize=2,
  showstringspaces=false,
  captionpos=b,
}

\lstdefinelanguage{Scala}{
  keywords={abstract,case,catch,class,def,do,else,extends,false,final,
           finally,for,forSome,if,implicit,import,lazy,match,new,null,
           object,override,package,private,protected,return,sealed,super,
           this,throw,trait,try,true,type,val,var,while,with,yield},
  ndkeywords={Boolean,Byte,Char,Class,Double,Float,Int,Long,Object,String,Unit},
  sensitive=true,
  morecomment=[l]{//},
  morecomment=[s]{/*}{*/},
  morestring=[b]",
  morestring=[b]',
}

% ============================================================
%  En-tetes / pieds de page
% ============================================================
\pagestyle{fancy}
\fancyhf{}
\fancyhead[L]{\small Fire Control System}
\fancyhead[R]{\small Modélisation et Vérification Formelle}
\fancyfoot[C]{\thepage}
\renewcommand{\headrulewidth}{0.4pt}
\renewcommand{\footrulewidth}{0pt}

% ============================================================
%  Style des titres
% ============================================================
\titleformat{\section}
  {\Large\bfseries\color{black!85}}
  {\thesection.}{0.6em}{}
\titleformat{\subsection}
  {\large\bfseries\color{black!75}}
  {\thesubsection}{0.6em}{}
\titleformat{\subsubsection}
  {\normalsize\bfseries\color{black!70}}
  {\thesubsubsection}{0.6em}{}

% ============================================================
%  Macros
% ============================================================
\newcommand{\fcs}{\textsc{Fcs}\xspace}
\newcommand{\code}[1]{\texttt{#1}}
\newcommand{\OK}{\textcolor{green!50!black}{\textbf{OK}}}
\newcommand{\KO}{\textcolor{red!70!black}{\textbf{KO}}}
\newcommand{\PASS}{\textcolor{green!50!black}{\textbf{PASS}}}

\newcolumntype{C}[1]{>{\centering\arraybackslash}p{#1}}
\newcolumntype{L}[1]{>{\raggedright\arraybackslash}p{#1}}

% ============================================================
%  Page de titre
% ============================================================
\title{
  \vspace{-1.5cm}
  \rule{\linewidth}{0.4pt}\\[0.4cm]
  {\Huge \bfseries Modélisation et Vérification Formelle\\[0.3em]
   d'un Système Distribué Critique}\\[0.4cm]
  {\Large \itshape Fire Control System (FCS)}\\[0.2cm]
  {\large Akka, Scala et Réseaux de Petri}\\[0.2cm]
  \rule{\linewidth}{0.4pt}\\[0.6cm]
  {\large Rapport détaillé du projet de Programmation Fonctionnelle}
}

\author{
  Paul Pitiot, Mathéo Crayssac \\
  \textit{ING 2 -- S2 -- CY Tech} \\[0.2em]
  \href{https://github.com/mcrayssac/Fire-Control-System}{\texttt{github.com/mcrayssac/Fire-Control-System}}
}

\date{Mai 2026}

% ============================================================
%  Document
% ============================================================
\begin{document}

\maketitle
\thispagestyle{empty}

\begin{abstract}
\noindent
Ce rapport présente la conception, la modélisation et la vérification formelle d'un
système distribué critique : le \emph{Fire Control System} (FCS), un système de
contrôle de tir embarqué impliquant détection, suivi balistique, gestion de munitions,
validation des règles d'engagement et exécution synchronisée. L'implémentation
s'appuie sur le modèle d'acteurs \textbf{Akka} en \textbf{Scala} et la traçabilité des
événements via \textbf{Kafka}. La vérification formelle est conduite à l'aide d'un
\textbf{réseau de Petri} dont l'analyseur est entièrement développé sans outil tiers,
conformément à la consigne. L'analyse couvre l'exploration exhaustive de l'espace
d'états (111 marquages, 227 arcs, \emph{0 deadlock}), le calcul des invariants de
places et de transitions par élimination de Gauss, la bornitude, la vivacité, ainsi
que la vérification de \textbf{10 invariants métier} et de \textbf{6 propriétés LTL}
de sûreté et de vivacité. Une simulation comparée valide la cohérence entre le
modèle formel et l'implémentation Akka sur six scénarios critiques. Le système
inclut également un mode interactif permettant de piloter la simulation pas à pas.
\end{abstract}

\vspace{0.3cm}
\begin{center}
\fbox{\parbox{0.85\linewidth}{\centering
\textbf{Mots-clés} : \emph{vérification formelle}, \emph{réseaux de Petri},
\emph{model checking}, \emph{LTL}, \emph{Akka}, \emph{Scala}, \emph{systèmes
distribués critiques}, \emph{actors}, \emph{Kafka}.
}}
\end{center}

\newpage
\tableofcontents
\newpage

% ============================================================
\section{Contexte, motivation et objectifs}
% ============================================================

\subsection{Contexte et choix applicatif}

Les systèmes distribués critiques (avionique, défense, industrie nucléaire,
finance) imposent un niveau d'exigence très élevé en matière de fiabilité.
Une simple violation d'invariant peut entraîner des conséquences irréversibles
(catastrophe humaine, perte matérielle majeure, atteinte à la sécurité publique).
L'échec du vol inaugural d'\emph{Ariane 5} en 1996, causé par un dépassement
d'entier non détecté lors d'une conversion 64\,bits vers 16\,bits, illustre l'utilité
des approches formelles couplées au test logiciel~\cite{lions1996ariane}.

L'application choisie pour ce projet est un \textbf{Fire Control System} (FCS),
système de contrôle de tir embarqué. Son choix se justifie par la conjonction des
caractéristiques suivantes :

\begin{itemize}[leftmargin=1.2cm,itemsep=2pt]
  \item \textbf{Criticité} : une erreur dans la séquence de tir produit des
        conséquences irréversibles ;
  \item \textbf{Concurrence} : capteur, suivi de cible, munitions et autorisation
        opèrent en parallèle ;
  \item \textbf{Synchronisation} : le tir n'est exécuté que lorsque toutes les
        préconditions sont simultanément réunies ;
  \item \textbf{Ressources finies} : le stock de munitions est explicitement modélisé ;
  \item \textbf{Tolérance aux pannes} : le système doit se rétablir d'une erreur
        sans rester bloqué.
\end{itemize}

\subsection{Objectifs du projet}

Conformément au sujet, le travail consiste à :
\begin{enumerate}[leftmargin=1.2cm,itemsep=2pt]
  \item établir un état de l'art rigoureux sur la vérification formelle, les
        réseaux de Petri et la logique LTL ;
  \item modéliser fonctionnellement et de manière concurrente le \fcs sous forme
        d'acteurs Akka, en identifiant les flux critiques ;
  \item traduire le modèle Akka en réseau de Petri, capturant l'espace d'états ;
  \item vérifier les propriétés structurelles (validité, absence de deadlock,
        cohérence) et les invariants métier ;
  \item simuler le comportement réel et le comparer aux prédictions formelles ;
  \item mettre à disposition l'ensemble du code source ouvert sur GitHub.
\end{enumerate}

\subsection{Note méthodologique}

Conformément à la remarque du sujet, \emph{aucun outil logiciel de réseaux de
Petri n'a été utilisé}. L'ensemble du moteur de réseau de Petri (franchissement,
exploration, P/T-invariants par élimination de Gauss, vérification LTL) est
implémenté en Scala par les auteurs, dans le package \code{fcs.petri}.

% ============================================================
\section{État de l'art}
% ============================================================

\subsection{Vérification formelle des systèmes critiques}

La \textbf{vérification formelle} consiste à prouver mathématiquement qu'un
système respecte certaines propriétés. Contrairement au test logiciel, qui
explore un nombre fini de scénarios, la vérification formelle analyse
\emph{l'ensemble} des comportements possibles.

\paragraph{Model checking vs.\ \textit{theorem proving}.}
Deux familles d'approches dominent :

\begin{table}[H]
\centering\small
\begin{tabularx}{\linewidth}{l X X X}
\toprule
\rowcolor{tableheader}\textbf{Approche} & \textbf{Principe} & \textbf{Avantages} & \textbf{Limites} \\
\midrule
Model checking
& Exploration exhaustive d'un modèle fini
& Automatique, fournit un contre-exemple
& Explosion combinatoire \\
\addlinespace
Theorem proving
& Déduction logique (Coq, Isabelle, Lean)
& Espace d'états infini possible
& Expertise élevée, intervention humaine \\
\bottomrule
\end{tabularx}
\caption{Comparaison des deux grandes familles de vérification formelle.}
\end{table}

Le \fcs comportant un espace d'états compact (111 marquages atteignables), le
\textbf{model checking} est l'approche retenue. Notre démarche consiste à :

\begin{enumerate}[leftmargin=1.2cm,itemsep=2pt]
  \item modéliser les états et transitions par un réseau de Petri ;
  \item explorer cet espace de façon exhaustive (BFS) ;
  \item vérifier 10 invariants métier et 6 propriétés LTL ;
  \item calculer les P/T-invariants par élimination de Gauss ;
  \item analyser bornitude et vivacité.
\end{enumerate}

\subsection{Réseaux de Petri}

\subsubsection{Définition formelle}

Un \textbf{réseau de Petri}~\cite{petri1962} est un quintuplet
\[
N \;=\; (P,\, T,\, \mathrm{Pre},\, \mathrm{Post},\, M_0)
\]
où :
\begin{itemize}[leftmargin=1.2cm,itemsep=2pt]
  \item $P$ est un ensemble fini de \textbf{places} (cercles) ;
  \item $T$ est un ensemble fini de \textbf{transitions} (barres) ;
  \item $\mathrm{Pre} : T \times P \to \mathbb{N}$ est la fonction de
        \emph{pré-condition} (arcs entrants) ;
  \item $\mathrm{Post} : T \times P \to \mathbb{N}$ est la fonction de
        \emph{post-condition} (arcs sortants) ;
  \item $M_0 : P \to \mathbb{N}$ est le \emph{marquage initial}.
\end{itemize}

Les \textbf{jetons} représentent l'état courant : un marquage $M$ associe à
chaque place son nombre de jetons.

\subsubsection{Sémantique de franchissement}

Une transition $t \in T$ est \textbf{franchissable} dans un marquage $M$
si et seulement si :
\[
\forall p \in P,\quad M(p) \;\geq\; \mathrm{Pre}(t, p).
\]
Le franchissement de $t$ produit un nouveau marquage $M'$ :
\[
\forall p \in P,\quad M'(p) \;=\; M(p) - \mathrm{Pre}(t, p) + \mathrm{Post}(t, p).
\]

\subsubsection{Matrice d'incidence et équation fondamentale}

La \textbf{matrice d'incidence} $C \in \mathbb{Z}^{|P| \times |T|}$ est définie par
\[
C \;=\; \mathrm{Post} - \mathrm{Pre}.
\]
L'équation fondamentale relie un marquage atteignable au marquage initial via
le \emph{vecteur de Parikh} $\sigma \in \mathbb{N}^{|T|}$ comptant les
franchissements :
\[
M' \;=\; M_0 + C \cdot \sigma.
\]

\subsubsection{Propriétés structurelles analysables}

\begin{table}[H]
\centering\small
\begin{tabularx}{\linewidth}{l X}
\toprule
\rowcolor{tableheader}\textbf{Propriété} & \textbf{Définition} \\
\midrule
Bornitude
& $N$ est $k$-borné si $\forall M$ atteignable, $\forall p \in P$,
  $M(p) \leq k$. Si $k=1$ on parle de réseau \emph{sauf}. \\
Vivacité
& Niveaux $L_0$ (morte), $L_1$ (potentiellement franchissable),
  $L_4$ (vivante depuis tout état atteignable). \\
Absence de deadlock
& Aucun état atteignable n'est dépourvu de transition franchissable. \\
P-invariant
& Vecteur $y \geq 0$ tel que $y^\top C = 0$ ; alors $y^\top M$ est constant
  (loi de conservation). \\
T-invariant
& Vecteur $x \geq 0$ tel que $C\,x = 0$ ; correspond à un comportement cyclique
  (effet net nul). \\
\bottomrule
\end{tabularx}
\caption{Propriétés structurelles standard d'un réseau de Petri.}
\end{table}

\subsection{Logique temporelle linéaire (LTL)}

Les invariants statiques ne suffisent pas pour exprimer toutes les propriétés
critiques d'un système réactif. La \textbf{logique temporelle linéaire}
(\emph{Linear Temporal Logic}, LTL), introduite par Pnueli en 1977~\cite{pnueli1977},
permet de raisonner sur l'évolution temporelle des états.

\paragraph{Opérateurs temporels.}
\begin{table}[H]
\centering\small
\begin{tabular}{l c l}
\toprule
\rowcolor{tableheader}\textbf{Opérateur} & \textbf{Notation} & \textbf{Signification} \\
\midrule
Globally
& $G(\varphi)$
& $\varphi$ est vraie dans \emph{tous} les états futurs \\
Finally
& $F(\varphi)$
& $\varphi$ sera vraie dans \emph{au moins un} état futur \\
Next
& $X(\varphi)$
& $\varphi$ est vraie dans l'état immédiatement suivant \\
Until
& $\varphi\, U\, \psi$
& $\varphi$ reste vraie jusqu'à ce que $\psi$ devienne vraie \\
\bottomrule
\end{tabular}
\caption{Opérateurs de la logique LTL.}
\end{table}

\paragraph{Types de propriétés.}
\begin{itemize}[leftmargin=1.2cm,itemsep=2pt]
  \item \textbf{Sûreté} (\emph{safety}) : « quelque chose de mauvais n'arrive
        jamais » $\rightarrow$ opérateur $G$.\\
        Exemple : $G(\neg(\textit{firing} \wedge \textit{reloading}))$.
  \item \textbf{Vivacité} (\emph{liveness}) : « quelque chose de bon finit par
        arriver » $\rightarrow$ opérateur $F$ ou $G(F(\varphi))$.\\
        Exemple : $G(\textit{error} \rightarrow F(\textit{idle}))$.
\end{itemize}

\paragraph{Vérification par model checking.}
\begin{itemize}[leftmargin=1.2cm,itemsep=2pt]
  \item $G(\varphi)$ : vérifier $\varphi$ dans chaque état atteignable ;
  \item $G(F(\varphi))$ : depuis chaque état, un état satisfaisant $\varphi$
        est toujours atteignable ;
  \item en cas de violation, on extrait un \emph{contre-exemple} (un chemin
        d'exécution fautif).
\end{itemize}

\subsection{Modèle d'acteurs et correspondance avec les réseaux de Petri}

Le \textbf{modèle d'acteurs}, formalisé par Hewitt en 1973~\cite{hewitt1973},
est un paradigme de calcul concurrent où chaque entité :
\begin{itemize}[leftmargin=1.2cm,itemsep=2pt]
  \item possède son propre état interne (pas de mémoire partagée) ;
  \item communique par envoi de messages asynchrones ;
  \item traite un seul message à la fois.
\end{itemize}

\textbf{Akka}~\cite{lightbend2024} est l'implémentation de ce modèle pour la
JVM. Akka Typed offre supervision hiérarchique et messages typés à la
compilation.

\paragraph{Correspondance Akka $\leftrightarrow$ Réseau de Petri.}
\begin{table}[H]
\centering\small
\begin{tabularx}{\linewidth}{l X}
\toprule
\rowcolor{tableheader}\textbf{Concept Akka} & \textbf{Concept Petri Net} \\
\midrule
État d'un acteur                       & Place contenant un jeton \\
Réception d'un message                 & Transition \\
Message entre acteurs                  & Arc \\
Ressource partagée (munitions)         & Place ressource (plusieurs jetons) \\
Synchronisation (attente de $N$ conditions) & Transition de synchronisation \\
\bottomrule
\end{tabularx}
\caption{Correspondance entre concepts Akka et concepts des réseaux de Petri.}
\end{table}

\noindent
L'intérêt de combiner les deux approches est complémentaire : le réseau de
Petri \emph{prouve} l'absence de deadlocks par exhaustivité, tandis que la
simulation Akka \emph{valide} le comportement réel par réalisme.

% ============================================================
\section{Architecture du système FCS}
% ============================================================

\subsection{Architecture en acteurs Akka}

L'application est organisée en \textbf{8 acteurs typés} placés sous une
hiérarchie de supervision :

\begin{lstlisting}[caption={Hiérarchie de supervision Akka.}]
SupervisorActor (racine)
|-- SensorActor          -- Detection de cibles
|-- TrackingActor        -- Calcul balistique
|-- AmmoActor            -- Gestion du stock de munitions
|-- CommandActor         -- Validation des ROE
|-- FireControlActor     -- Orchestrateur central
|-- KafkaProducerActor   -- Publication d'evenements
+-- KafkaConsumerActor   -- Audit trail
\end{lstlisting}

\begin{table}[H]
\centering\small
\begin{tabularx}{\linewidth}{l X}
\toprule
\rowcolor{tableheader}\textbf{Acteur} & \textbf{Responsabilité} \\
\midrule
\code{SupervisorActor}     & Création, supervision et redémarrage des acteurs enfants. \\
\code{SensorActor}         & Détection de cibles ; déclenchement du cycle de tir. \\
\code{TrackingActor}       & Calcul balistique (élévation, dérive, temps de vol, indice de confiance). \\
\code{AmmoActor}           & Gestion du stock de munitions (APFSDS, HEAT, HESH, HE), chargement. \\
\code{CommandActor}        & Validation des règles d'engagement (confiance min., portée max.). \\
\code{FireControlActor}    & Orchestrateur ; attend les 3 confirmations, exécute le tir, gère reload/cooldown. \\
\code{KafkaProducerActor}  & Publication sur les 7 topics Kafka. \\
\code{KafkaConsumerActor}  & Collecte de l'\emph{audit trail}. \\
\bottomrule
\end{tabularx}
\caption{Rôle de chaque acteur du \fcs.}
\end{table}

\subsection{Flux du cycle de tir nominal}

Le cycle suivant illustre l'interaction des acteurs lors d'un tir réussi :

\begin{lstlisting}[caption={Diagramme de séquence simplifié du cycle nominal.}]
SensorActor              FireControlActor             TrackingActor
    |                          |                            |
    |  SimulateDetection       |                            |
    |------------------------->|  TrackTarget               |
    |                          |--------------------------->|
    |                          |  TargetLockConfirmed       |
    |                          |<---------------------------|
    |                          |
    |                          |  LoadAmmo  ----------> AmmoActor
    |                          |  AmmoLoadConfirmed <------'
    |                          |
    |                          |  RequestAuth -------> CommandActor
    |                          |  FireAuthConfirmed <------'
    |                          |
    |                          |  [readyToFire = true]
    |                          |  executeFire()
    |                          |  PublishEvent ------> KafkaProducerActor
    |                          |
    |                          |  ReloadTimeout (2s)
    |                          |  CooldownTimeout (3s)
    |                          |  --> retour a Idle
\end{lstlisting}

\paragraph{Point de synchronisation.}
Le \code{FireControlActor} attend trois confirmations distinctes avant
d'autoriser le tir :
\begin{enumerate}[leftmargin=1.2cm,itemsep=1pt]
  \item \code{TargetLockConfirmed} -- verrouillage balistique réussi ;
  \item \code{AmmoLoadConfirmed}   -- munition chargée ;
  \item \code{FireAuthConfirmed}   -- règles d'engagement validées.
\end{enumerate}
La condition de tir s'écrit :
\[
\textit{readyToFire} \;=\; \textit{targetLocked} \,\wedge\, \textit{ammoLoaded}
\,\wedge\, \textit{fireAuthorized}.
\]
Cette logique de fusion correspond à la \textbf{transition $T_4$} du réseau de
Petri (synchronisation $P_3 + P_4 \rightarrow P_5$).

\subsection{Gestion des erreurs}

\begin{table}[H]
\centering\small
\begin{tabularx}{\linewidth}{l X}
\toprule
\rowcolor{tableheader}\textbf{Situation} & \textbf{Comportement} \\
\midrule
Échec verrouillage (confiance $<$ 0.5)
& \code{TrackingActor} renvoie \code{TargetLockFailed} ; cycle avorté. \\
Refus d'autorisation (confiance $<$ 0.8 ou portée $>$ 4000\,m)
& \code{CommandActor} renvoie \code{FireAuthDenied} ; cycle avorté. \\
Stock épuisé
& \code{AmmoActor} renvoie \code{AmmoLoadFailed} ; cycle avorté. \\
Erreur système
& \code{SupervisorActor} émet \code{ReportError} ; stratégie \emph{restart}. \\
\bottomrule
\end{tabularx}
\caption{Modes d'erreur gérés par l'architecture.}
\end{table}

\subsection{Topics Kafka}

L'audit trail repose sur six topics Kafka couvrant l'ensemble du cycle :

\begin{table}[H]
\centering\small
\begin{tabular}{l l l}
\toprule
\rowcolor{tableheader}\textbf{Topic} & \textbf{Événement} & \textbf{Producteur} \\
\midrule
\code{fcs.target.detected}  & Cible détectée       & \code{SensorActor} \\
\code{fcs.target.locked}    & Solution balistique  & \code{TrackingActor} \\
\code{fcs.fire.authorized}  & ROE validées         & \code{CommandActor} \\
\code{fcs.fire.executed}    & Tir effectué         & \code{FireControlActor} \\
\code{fcs.ammo.status}      & Munition chargée     & \code{AmmoActor} \\
\code{fcs.error.critical}   & Erreur fatale        & \code{SupervisorActor} \\
\bottomrule
\end{tabular}
\caption{Topics Kafka et leurs producteurs.}
\end{table}

% ============================================================
\section{Modèle formel : réseau de Petri du FCS}
% ============================================================

\subsection{Définition du réseau}

Le réseau de Petri du \fcs est défini par
$N = (P, T, \mathrm{Pre}, \mathrm{Post}, M_0)$ avec
\[
|P| = 13 \quad\text{et}\quad |T| = 12.
\]
Le marquage initial est
\[
M_0 \;=\; (1, 0, 0, 0, 0, 0, 0, 0, 0, 0, n, 0, 0)
\]
où $n$ désigne le stock initial de munitions (typiquement $n = 3$ dans les
expériences).

\subsubsection{Places}

\begin{table}[H]
\centering\small
\begin{tabular}{c l l l}
\toprule
\rowcolor{tableheader}\textbf{Place} & \textbf{Nom} & \textbf{Type} & \textbf{Rôle} \\
\midrule
$P_0$  & \code{Idle}            & Contrôle  & Système au repos \\
$P_1$  & \code{Target\_Detected}& Contrôle  & Cible détectée \\
$P_2$  & \code{Target\_Locked}  & Contrôle  & Verrouillage confirmé \\
$P_3$  & \code{Ammo\_Loaded}    & Contrôle  & Munition chargée \\
$P_4$  & \code{Fire\_Authorized}& Contrôle  & Autorisation accordée \\
$P_5$  & \code{Ready\_To\_Fire} & Contrôle  & Préconditions réunies \\
$P_6$  & \code{Firing}          & Contrôle  & Tir en cours \\
$P_7$  & \code{Reloading}       & Contrôle  & Rechargement \\
$P_8$  & \code{Cooldown}        & Contrôle  & Refroidissement \\
$P_9$  & \code{Error\_State}    & Contrôle  & Erreur détectée \\
$P_{10}$ & \code{Ammo\_Stock}   & Ressource & Stock de munitions \\
$P_{11}$ & \code{Kafka\_Queue}  & Ressource & Événements en attente \\
$P_{12}$ & \code{Log\_Recorded} & Ressource & Événements journalisés \\
\bottomrule
\end{tabular}
\caption{Places du réseau de Petri du \fcs.}
\end{table}

\subsubsection{Transitions}

\begin{table}[H]
\centering\small
\begin{tabular}{c l l}
\toprule
\rowcolor{tableheader}\textbf{Transition} & \textbf{Nom} & \textbf{$\mathrm{Pre} \to \mathrm{Post}$} \\
\midrule
$T_0$    & \code{detect\_target}     & $\{P_0, P_{10}\} \to \{P_1, P_{10}, P_{11}\}$ \\
$T_1$    & \code{lock\_target}       & $\{P_1\} \to \{P_2\}$ \\
$T_2$    & \code{load\_ammo}         & $\{P_{10}\} \to \{P_3\}$ \\
$T_3$    & \code{authorize\_fire}    & $\{P_2\} \to \{P_4\}$ \\
$T_4$    & \code{ready\_sync}        & $\{P_3, P_4\} \to \{P_5\}$ \\
$T_5$    & \code{fire}               & $\{P_5\} \to \{P_6, P_{11}\}$ \\
$T_6$    & \code{end\_fire}          & $\{P_6\} \to \{P_7\}$ \\
$T_7$    & \code{reload\_complete}   & $\{P_7\} \to \{P_8\}$ \\
$T_8$    & \code{cooldown\_complete} & $\{P_8\} \to \{P_0\}$ \\
$T_9$    & \code{error\_detected}    & $\{P_0\} \to \{P_9\}$ \\
$T_{10}$ & \code{error\_recovery}    & $\{P_9\} \to \{P_0\}$ \\
$T_{11}$ & \code{kafka\_log}         & $\{P_{11}\} \to \{P_{12}\}$ \\
\bottomrule
\end{tabular}
\caption{Transitions du réseau de Petri du \fcs.}
\end{table}

\subsection{Correspondance Acteur Akka $\leftrightarrow$ Réseau de Petri}

\begin{table}[H]
\centering\small
\begin{tabularx}{\linewidth}{l X X}
\toprule
\rowcolor{tableheader}\textbf{Acteur Akka} & \textbf{Places} & \textbf{Transitions} \\
\midrule
\code{SensorActor}          & $P_0,\ P_1$
                            & $T_0$ \\
\code{TrackingActor}        & $P_1 \to P_2$
                            & $T_1$ \\
\code{AmmoActor}            & $P_{10},\ P_3$
                            & $T_2$ \\
\code{CommandActor}         & $P_2 \to P_4$
                            & $T_3$ \\
\code{FireControlActor}     & $P_3 + P_4 \to P_5 \to P_6 \to P_7 \to P_8 \to P_0$
                            & $T_4, T_5, T_6, T_7, T_8$ \\
\code{SupervisorActor}      & $P_0 \leftrightarrow P_9$
                            & $T_9, T_{10}$ \\
\code{KafkaProducerActor}   & $P_{11} \to P_{12}$
                            & $T_{11}$ \\
\bottomrule
\end{tabularx}
\caption{Correspondance entre acteurs Akka et éléments du réseau.}
\end{table}

\subsection{Cycle nominal pas à pas}

\begin{table}[H]
\centering\small
\begin{tabular}{c c l l l}
\toprule
\rowcolor{tableheader}\textbf{\#} & \textbf{Trans.} & \textbf{Acteur} & \textbf{Message Akka} & \textbf{Effet Petri} \\
\midrule
1  & $T_0$    & \code{SensorActor}        & \code{SimulateDetection}      & $P_0,P_{10} \to P_1,P_{10},P_{11}$ \\
2  & $T_1$    & \code{TrackingActor}      & \code{TargetLockConfirmed}    & $P_1 \to P_2$ \\
3  & $T_2$    & \code{AmmoActor}          & \code{AmmoLoadConfirmed}      & $P_{10} \to P_3$ \\
4  & $T_3$    & \code{CommandActor}       & \code{FireAuthConfirmed}      & $P_2 \to P_4$ \\
5  & $T_4$    & \code{FireControlActor}   & \code{readyToFire = true}     & $P_3,P_4 \to P_5$ \\
6  & $T_5$    & \code{FireControlActor}   & \code{executeFire()}          & $P_5 \to P_6,P_{11}$ \\
7  & $T_6$    & \code{FireControlActor}   & (interne)                     & $P_6 \to P_7$ \\
8  & $T_7$    & \code{FireControlActor}   & \code{ReloadTimeout}          & $P_7 \to P_8$ \\
9  & $T_8$    & \code{FireControlActor}   & \code{CooldownTimeout}        & $P_8 \to P_0$ \\
10 & $T_{11}$ & \code{KafkaProducerActor} & \code{PublishEvent} ($\times2$)& $P_{11} \to P_{12}$ \\
\bottomrule
\end{tabular}
\caption{Cycle nominal de tir, transition par transition.}
\end{table}

Les étapes 2, 3 et 4 sont \emph{réellement concurrentes} dans Akka (messages
asynchrones). Dans le réseau de Petri, leur entrelacement quelconque produit
le même état post-synchronisation grâce à la transition de fusion $T_4$.

\subsection{Matrice d'incidence}

La matrice $C = \mathrm{Post} - \mathrm{Pre}$, de dimensions $13 \times 12$,
est explicitée ci-dessous (les zéros sont représentés par « $\cdot$ ») :

\begin{lstlisting}[basicstyle=\ttfamily\footnotesize,caption={Matrice d'incidence du réseau de Petri du \fcs.}]
                   T0   T1   T2   T3   T4   T5   T6   T7   T8   T9   T10  T11
Idle               -1    .    .    .    .    .    .    .   +1   -1   +1    .
Target_Detected    +1   -1    .    .    .    .    .    .    .    .    .    .
Target_Locked       .   +1    .   -1    .    .    .    .    .    .    .    .
Ammo_Loaded         .    .   +1    .   -1    .    .    .    .    .    .    .
Fire_Authorized     .    .    .   +1   -1    .    .    .    .    .    .    .
Ready_To_Fire       .    .    .    .   +1   -1    .    .    .    .    .    .
Firing              .    .    .    .    .   +1   -1    .    .    .    .    .
Reloading           .    .    .    .    .    .   +1   -1    .    .    .    .
Cooldown            .    .    .    .    .    .    .   +1   -1    .    .    .
Error_State         .    .    .    .    .    .    .    .    .   +1   -1    .
Ammo_Stock          .    .   -1    .    .    .    .    .    .    .    .    .
Kafka_Queue        +1    .    .    .    .   +1    .    .    .    .    .   -1
Log_Recorded        .    .    .    .    .    .    .    .    .    .    .   +1
\end{lstlisting}

\paragraph{Lecture.} La colonne $T_2$ indique que la place $P_{10}$
(\code{Ammo\_Stock}) perd un jeton ($-1$) et que $P_3$ (\code{Ammo\_Loaded})
en gagne un ($+1$) : c'est la consommation d'une munition.

\subsection{Éléments non modélisés}

Le modèle formel abstrait certains détails techniques jugés non pertinents
pour la sûreté. Ces choix sont motivés par le souci de garder l'espace
d'états compact tout en préservant les propriétés étudiées.

\begin{table}[H]
\centering\small
\begin{tabularx}{\linewidth}{l X}
\toprule
\rowcolor{tableheader}\textbf{Aspect Akka} & \textbf{Raison de l'abstraction} \\
\midrule
Types de munitions (APFSDS, HEAT, HESH, HE)
& Stock générique ($P_{10}$) ; aucun impact sur la sûreté de la séquence. \\
Calcul balistique détaillé
& Abstrait en $T_1$ (succès / échec). \\
Timeouts (tracking 5s, reload 2s, cooldown 3s)
& Réseau \emph{non temporisé} ; les ordres causals suffisent. \\
Scénarios multiples du Supervisor
& Cycle générique $T_9 \to T_{10}$. \\
Supervision et restart Akka
& Abstraits par les transitions d'erreur et de récupération. \\
\bottomrule
\end{tabularx}
\caption{Détails techniques non modélisés et justification.}
\end{table}

Le diagramme complet est fourni dans \code{docs/fcs\_petri\_net.drawio}.

% ============================================================
\section{Vérification des propriétés}
% ============================================================

\subsection{Espace d'états}

L'analyseur explore exhaustivement l'espace d'états par BFS à partir de
$M_0 = (P_0=1, P_{10}=3)$ :

\begin{table}[H]
\centering\small
\begin{tabular}{l c}
\toprule
\rowcolor{tableheader}\textbf{Métrique} & \textbf{Valeur} \\
\midrule
Marquages atteignables  & \textbf{111} \\
Arcs                    & \textbf{227} \\
Deadlocks               & \textbf{0}   \\
\bottomrule
\end{tabular}
\caption{Synthèse de l'exploration de l'espace d'états.}
\end{table}

L'absence de deadlock signifie que depuis tout marquage atteignable,
au moins une transition reste franchissable.

\subsection{Analyse structurelle}

\subsubsection{P-invariants}

Par élimination de Gauss sur le système $C^\top y = 0$, on obtient
\textbf{deux P-invariants minimaux} :

\paragraph{$y_1$ -- conservation du flux de contrôle.}
\[
y_1 = (1,1,1,0,1,1,1,1,1,1,0,0,0)
\]
\[
P_0 + P_1 + P_2 + P_4 + P_5 + P_6 + P_7 + P_8 + P_9 \;=\; 1.
\]
Le jeton de contrôle est toujours dans \emph{exactement une} place du cycle
principal. Cet invariant correspond à l'invariant métier \textbf{INV7}.

\paragraph{$y_2$ -- conservation globale.}
\[
y_2 = (1,0,0,2,0,2,1,1,1,1,2,1,1)
\]
relie les places de contrôle aux places ressource. La somme pondérée totale
reste constante au cours de toute exécution.

Les deux P-invariants sont vérifiés sur les 111 marquages atteignables.

\subsubsection{T-invariants}

La résolution de $C\,x = 0$ identifie \textbf{un T-invariant} minimal :
\[
x_1 = (0,0,0,0,0,0,0,0,0,1,1,0) \;=\; \{T_9, T_{10}\}.
\]
Le cycle erreur / récupération est cyclique au sens des T-invariants : il
restaure le marquage initial. En revanche, le cycle nominal de tir
\emph{n'est pas} un T-invariant : il consomme une munition et produit deux
événements journalisés (effet net non nul, ce qui est la marque d'un
processus irréversible).

\subsection{Bornitude}

\begin{table}[H]
\centering\small
\begin{tabular}{l c l}
\toprule
\rowcolor{tableheader}\textbf{Place(s)} & \textbf{Borne max} & \textbf{Commentaire} \\
\midrule
$P_0$--$P_2$, $P_4$--$P_9$ (contrôle) & 1 & Toutes 1-bornées (\emph{sauf}). \\
$P_3$ (\code{Ammo\_Loaded})           & 3 & Bornée par le stock initial. \\
$P_{10}$ (\code{Ammo\_Stock})         & 3 & Stock initial. \\
$P_{11}$ (\code{Kafka\_Queue})        & 2 & Au plus 2 événements en file. \\
$P_{12}$ (\code{Log\_Recorded})       & 6 & 2 événements par cycle $\times$ 3 cycles max. \\
\bottomrule
\end{tabular}
\caption{Bornitude par place.}
\end{table}

\noindent
\textbf{Résultat global :} le réseau est \textbf{6-borné} ($k=6$, dû à
$P_{12}$). Toutes les places de contrôle sont 1-bornées, ce qui constitue
une garantie de \emph{stabilité} du flux de contrôle.

\subsection{Vivacité}

\begin{table}[H]
\centering\small
\begin{tabular}{l c l}
\toprule
\rowcolor{tableheader}\textbf{Transition} & \textbf{Niveau} & \textbf{Commentaire} \\
\midrule
$T_0$--$T_8$, $T_{11}$  & $L_1$ & Potentiellement franchissable. \\
$T_9, T_{10}$           & $L_4$ & Vivantes (cycle erreur réinitialisable). \\
\bottomrule
\end{tabular}
\caption{Niveaux de vivacité par transition.}
\end{table}

\noindent
Le réseau \emph{n'est pas} vivant au sens $L_4$ \emph{dans son ensemble} :
après épuisement des munitions ($P_{10}=0$), les transitions du cycle de tir
ne sont plus franchissables. Seul le cycle erreur ($T_9, T_{10}$) reste
vivant. Ce comportement est \textbf{physiquement correct} : un système sans
munitions ne peut plus tirer. Aucune transition n'est morte (\,$L_0$\,).

\subsection{Invariants métier}

Les invariants métier traduisent les contraintes opérationnelles non
négociables du \fcs : règles d'engagement, conservation des ressources,
exclusions mutuelles entre phases du cycle de tir et propriétés de
\emph{liveness} attendues du système. Pour chaque invariant nous présentons
successivement sa \emph{motivation métier}, sa \emph{formalisation} sur le
réseau de Petri, la \emph{méthode de vérification} appliquée par
l'analyseur, et l'\emph{interprétation} du résultat obtenu sur les 111
marquages atteignables. La synthèse est rappelée dans le tableau
récapitulatif en fin de section.

% ----------- INV1 -----------
\subsubsection*{INV1 -- Pas de tir sans verrouillage de cible}

\paragraph{Motivation.}
Tirer sans solution balistique consolidée revient à projeter une munition
sur une trajectoire incertaine : c'est la cause type d'un \emph{tir fratricide}
ou d'un dommage collatéral. La règle métier exige que toute exécution
de la transition de tir soit \emph{causalement} précédée du verrouillage
($T_1$) puis de l'autorisation ($T_3$), elle-même conditionnée par le
verrouillage.

\paragraph{Formalisation.}
\[
\forall M \in \mathrm{Reach}(M_0),\quad
\bigl[T_5\ \text{franchissable en}\ M\bigr] \;\Longrightarrow\; \exists\, M' \prec M,\ M'(P_2) \geq 1.
\]
Sur le réseau, $T_5$ requiert $P_5 \geq 1$ ; or $P_5$ ne peut être marquée
que par $T_4$, qui consomme $P_4$ ; et $P_4$ n'est produit que par $T_3$
qui consomme $P_2$. La chaîne de pré-conditions est donc structurellement
forcée.

\paragraph{Méthode de vérification.}
Analyse par accessibilité des transitions : pour chaque marquage atteignable
où $T_5$ est franchissable, on vérifie qu'il existe au moins un prédécesseur
ayant $M(P_2) \geq 1$. La preuve repose en plus sur le P-invariant $y_1$
qui exclut tout marquage où $P_5$ serait marquée sans passage par $P_2$.

\paragraph{Résultat.}
\OK\ -- vérifié sur les 111 marquages. Aucun chemin du graphe d'accessibilité
ne mène à $T_5$ sans avoir auparavant traversé $T_1$.

% ----------- INV2 -----------
\subsubsection*{INV2 -- Pas de tir sans munition chargée}

\paragraph{Motivation.}
Il s'agit de garantir l'absence de \emph{tir à vide} : ni perte d'énergie
mécanique, ni faux positif sur l'audit Kafka. Pour un système réel, un tir
sans munition entraîne par ailleurs une usure inutile du mécanisme et un
faux signal vers la chaîne de commandement.

\paragraph{Formalisation.}
\[
\forall M \in \mathrm{Reach}(M_0),\quad T_5\ \text{franch.}\ \Longrightarrow\ \exists\, M' \prec M,\ M'(P_3) \geq 1.
\]
$T_5$ exige $M(P_5) \geq 1$ ; or $T_4$ consomme conjointement $P_3$ et
$P_4$. La place $P_3$ ne peut elle-même être produite que par $T_2$, qui
consomme un jeton de $P_{10}$.

\paragraph{Méthode de vérification.}
On démontre par induction sur la profondeur du BFS que tout marquage
satisfaisant $M(P_5) \geq 1$ a été précédé d'un marquage où $P_3 \geq 1$.
La transition de synchronisation $T_4$ joue ici un rôle structurant : elle
\emph{absorbe} simultanément le jeton de chargement et celui d'autorisation.

\paragraph{Résultat.}
\OK\ -- vérifié. La transition $T_5$ n'est jamais franchissable sans qu'une
munition n'ait été préalablement chargée par l'\code{AmmoActor}.

% ----------- INV3 -----------
\subsubsection*{INV3 -- Pas de tir sans autorisation des règles d'engagement}

\paragraph{Motivation.}
La conformité juridique et opérationnelle du tir dépend de l'aval explicite
du \code{CommandActor}, qui valide les \emph{Rules of Engagement} (ROE) :
seuil de confiance, portée, identification ami/ennemi. Cet invariant est la
contrepartie formelle de la chaîne de commandement.

\paragraph{Formalisation.}
\[
\forall M \in \mathrm{Reach}(M_0),\quad T_5\ \text{franch.}\ \Longrightarrow\ \exists\, M' \prec M,\ M'(P_4) \geq 1.
\]

\paragraph{Méthode de vérification.}
Identique à INV1/INV2 : analyse de la garde de $T_4$ qui exige $P_3 \wedge P_4$.
Le $P$-invariant $y_1$ confirme par ailleurs qu'un seul jeton de contrôle
circule, ce qui empêche la coexistence de deux flux distincts qui
contourneraient la garde.

\paragraph{Résultat.}
\OK\ -- vérifié. Le scénario 2 de la simulation comparée (\S5.2) en fournit
un test \emph{positif} : sans franchissement de $T_3$, $T_4$ se bloque, et
le cycle est avorté à l'identique côté Akka via \code{FireAuthDenied}.

% ----------- INV4 -----------
\subsubsection*{INV4 -- Stock de munitions toujours positif ou nul}

\paragraph{Motivation.}
Empêcher tout \emph{underflow} du stock physique : on ne peut consommer
qu'une munition réellement présente. Un stock négatif équivaudrait à un bug
classique d'imputation comptable -- ici dans un contexte ressources
critiques.

\paragraph{Formalisation.}
\[
\forall M \in \mathrm{Reach}(M_0),\quad M(P_{10}) \;\geq\; 0.
\]

\paragraph{Méthode de vérification.}
La sémantique de franchissement d'un réseau de Petri impose
$M(p) \geq \mathrm{Pre}(t,p)$ pour qu'une transition soit franchissable.
La place $P_{10}$ est consommée uniquement par $T_2$ (et lue par $T_0$ qui
la restitue), donc tout franchissement préserve $M(P_{10}) \geq 0$ par
construction. Vérification en outre par \emph{model checking} sur les 111
marquages : aucun n'affiche $M(P_{10}) < 0$.

\paragraph{Résultat.}
\OK\ -- vérifié. La borne maximale observée pour $P_{10}$ est égale au
stock initial $n=3$ ; la borne minimale est $0$, atteinte après trois
cycles complets de tir.

% ----------- INV5 -----------
\subsubsection*{INV5 -- Exclusion mutuelle entre tir et rechargement}

\paragraph{Motivation.}
Tirer pendant un rechargement est physiquement impossible et logiquement
dangereux : il s'agit d'une \emph{race condition} typique sur l'organe de
chargement. L'invariant garantit que les deux phases ne se chevauchent
jamais.

\paragraph{Formalisation.}
\[
\forall M \in \mathrm{Reach}(M_0),\quad M(P_6) + M(P_7) \;\leq\; 1.
\]

\paragraph{Méthode de vérification.}
Cette borne se déduit du $P$-invariant $y_1$ qui contraint la somme
$P_0 + P_1 + P_2 + P_4 + P_5 + P_6 + P_7 + P_8 + P_9 = 1$. Elle est
également vérifiée par énumération sur les 111 marquages atteignables.
Aucun marquage ne présente simultanément $P_6 = 1$ et $P_7 = 1$.

\paragraph{Résultat.}
\OK\ -- vérifié structurellement (par le P-invariant) et empiriquement
(par énumération exhaustive).

% ----------- INV6 -----------
\subsubsection*{INV6 -- Pas de tir pendant le cooldown}

\paragraph{Motivation.}
Le \emph{cooldown} thermique du canon est nécessaire pour éviter
l'endommagement matériel et les déformations du tube qui altèreraient la
balistique des tirs ultérieurs. Tirer pendant cette phase invaliderait la
règle de protection matérielle.

\paragraph{Formalisation.}
\[
\forall M \in \mathrm{Reach}(M_0),\quad M(P_6) > 0 \;\Longrightarrow\; M(P_8) = 0.
\]
De manière équivalente : $M(P_6) \cdot M(P_8) = 0$ partout.

\paragraph{Méthode de vérification.}
Le P-invariant $y_1$ implique que la somme des places de contrôle vaut $1$,
donc la coexistence $P_6 \wedge P_8$ est structurellement interdite.
Vérification empirique : sur les 111 marquages, l'analyseur teste
$M(P_6) > 0 \Rightarrow M(P_8) = 0$.

\paragraph{Résultat.}
\OK\ -- vérifié. Aucun marquage atteignable n'a simultanément $P_6 = 1$
et $P_8 = 1$.

% ----------- INV7 -----------
\subsubsection*{INV7 -- Conservation du flux de contrôle}

\paragraph{Motivation.}
Garantir l'\emph{unicité du jeton de contrôle} : à tout instant, le système
se trouve dans \emph{exactement un} état du cycle principal. Cette unicité
est au cœur de la déterminabilité du système : elle empêche les états
fantômes (deux activités simultanées) et les états orphelins (aucune
activité).

\paragraph{Formalisation.}
\[
\forall M \in \mathrm{Reach}(M_0),\quad
\sum_{p \in \{P_0,P_1,P_2,P_4,P_5,P_6,P_7,P_8,P_9\}} M(p) \;=\; 1.
\]
Soit, en notation vectorielle, $y_1^\top \cdot M = 1$ avec
$y_1 = (1,1,1,0,1,1,1,1,1,1,0,0,0)$.

\paragraph{Méthode de vérification.}
\begin{enumerate}[leftmargin=1.2cm,itemsep=1pt]
  \item Calcul algébrique : on vérifie que $y_1^\top \cdot C = 0$ par
        produit matriciel sur la matrice d'incidence (élimination de
        Gauss).
  \item Évaluation directe sur $M_0$ : $y_1^\top \cdot M_0 = 1$.
  \item Conséquence : $y_1^\top \cdot M = y_1^\top \cdot M_0 = 1$ pour tout
        marquage atteignable.
\end{enumerate}

\paragraph{Résultat.}
\OK\ -- vérifié algébriquement et empiriquement. INV7 fournit en outre la
preuve structurelle de INV5 et INV6.

% ----------- INV8 -----------
\subsubsection*{INV8 -- Tout tir est journalisé}

\paragraph{Motivation.}
La traçabilité est une exigence forte des systèmes critiques : tout tir
doit laisser une trace persistante exploitable a posteriori (audit,
enquête, après-action). Cet invariant relève de la \emph{vivacité} -- il ne
suffit pas que la trace puisse exister, elle doit \emph{finir par} exister.

\paragraph{Formalisation.}
\[
G\bigl(\textit{firing} \;\Longrightarrow\; F(\textit{log\_recorded})\bigr).
\]
À chaque entrée dans la place $P_6$, un jeton est ajouté à $P_{11}$
(par $T_5$), et $T_{11}$ déplace ce jeton vers $P_{12}$.

\paragraph{Méthode de vérification.}
Vérification par \emph{model checking} LTL : pour chaque marquage
satisfaisant $\textit{firing}$, on vérifie qu'au moins un chemin atteignable
mène à un état satisfaisant $\textit{log\_recorded}$. L'analyseur explore
récursivement les successeurs et garantit qu'aucun cycle ne piège l'exécution.

\paragraph{Résultat.}
\OK\ -- vérifié. La transition $T_{11}$ est toujours franchissable lorsque
$P_{11} \geq 1$, ce qui est garanti après chaque franchissement de $T_5$.

% ----------- INV9 -----------
\subsubsection*{INV9 -- Le système peut toujours retourner à l'état \textit{Idle}}

\paragraph{Motivation.}
Une propriété de \emph{réinitialisation} indispensable pour la maintenance
opérationnelle : le système ne doit jamais s'enliser dans un état
intermédiaire ; il doit toujours pouvoir, par une suite finie de
transitions, revenir à $P_0$ et être prêt pour un nouveau cycle.

\paragraph{Formalisation.}
\[
\forall M \in \mathrm{Reach}(M_0),\ \exists\, M' \in \mathrm{Reach}(M),\ M'(P_0) \geq 1.
\]

\paragraph{Méthode de vérification.}
Recherche dans le graphe d'accessibilité. Pour chaque marquage atteignable,
l'analyseur cherche par BFS un descendant satisfaisant $M(P_0) \geq 1$. Ce
test est positif pour les 111 marquages.

\paragraph{Résultat.}
\OK\ -- vérifié. Plus précisément, le quotient
$|\{M : M(P_0) \geq 1\}| / |\mathrm{Reach}(M_0)|$ est non nul, et chaque
état mène toujours à un tel marquage en au plus 4 franchissements.

% ----------- INV10 -----------
\subsubsection*{INV10 -- Absence de deadlock}

\paragraph{Motivation.}
Le \emph{deadlock} est l'erreur la plus redoutée des systèmes concurrents :
un état figé d'où aucune transition ne peut plus se déclencher. Pour un
système d'armement, un deadlock signifierait l'impossibilité de tirer, de
recharger, ou même de retourner à l'état de repos.

\paragraph{Formalisation.}
\[
\forall M \in \mathrm{Reach}(M_0),\ \exists\, t \in T,\ t\ \text{est franchissable en}\ M.
\]

\paragraph{Méthode de vérification.}
Pendant l'exploration BFS, l'analyseur compte pour chaque marquage le
nombre de transitions franchissables (\code{enabledTransitions}). Tout
marquage présentant $|\textit{enabled}(M)| = 0$ est marqué comme deadlock.
L'analyseur produit la liste des deadlocks ; ici, cette liste est vide.

\paragraph{Résultat.}
\OK\ -- 0 deadlock sur 111 marquages atteignables. Cette propriété est
\emph{globale} et nécessite l'exploration exhaustive : aucune approche
locale (statique) ne peut la garantir.

% ----------- Synthèse INV -----------
\subsubsection*{Synthèse des invariants métier}

\begin{table}[H]
\centering\small
\begin{tabularx}{\linewidth}{c X l c}
\toprule
\rowcolor{tableheader}\textbf{\#} & \textbf{Propriété} & \textbf{Formalisation} & \textbf{Résultat} \\
\midrule
INV1  & Pas de tir sans verrouillage      & $T_5$ franch. $\Rightarrow$ trace passe par $T_1$ & \OK \\
INV2  & Pas de tir sans munition chargée  & $T_4$ exige $P_3 \geq 1$                          & \OK \\
INV3  & Pas de tir sans autorisation      & $T_4$ exige $P_4 \geq 1$                          & \OK \\
INV4  & Stock munitions $\geq 0$          & $\forall M,\ M(P_{10}) \geq 0$                    & \OK \\
INV5  & Exclusion mutuelle tir/recharge   & $\forall M,\ M(P_6) + M(P_7) \leq 1$              & \OK \\
INV6  & Pas de tir pendant cooldown       & $M(P_6) > 0 \Rightarrow M(P_8) = 0$               & \OK \\
INV7  & Conservation du flux de contrôle  & $y_1^\top \cdot M = 1$                            & \OK \\
INV8  & Tout tir est journalisé           & $G(\textit{firing} \rightarrow F(\textit{log}))$  & \OK \\
INV9  & Retour à l'état \textit{idle}     & $\forall M,\ \exists M' \succeq M,\ P_0 \geq 1$   & \OK \\
INV10 & Absence de deadlock               & $\forall M, \exists t \in T,\ t\ \text{franch.}$  & \OK \\
\bottomrule
\end{tabularx}
\caption{Synthèse de la vérification des 10 invariants métier.}
\end{table}

\noindent
\textbf{Bilan : 10/10 invariants satisfaits.} Les invariants INV4--INV7 sont
prouvés \emph{structurellement} (P-invariants, sémantique des arcs), tandis
que INV1--INV3, INV8--INV10 reposent sur l'\emph{exploration exhaustive} du
graphe d'accessibilité. Cette double approche -- algébrique et énumérative
-- offre la robustesse caractéristique du \emph{model checking}.

\subsection{Propriétés LTL}

Les propriétés LTL ciblent l'\emph{évolution temporelle} du système, là où
les invariants métier ne décrivent que des états instantanés. Nous
distinguons trois familles : la \textbf{sûreté} ($G(\varphi)$ -- une mauvaise
chose ne se produit jamais), la \textbf{vivacité} ($F(\varphi)$ ou
$G(F(\varphi))$ -- une bonne chose finit par arriver) et la
\textbf{récupération} ($G(\varphi \rightarrow F(\psi))$ -- depuis tout état
$\varphi$, on revient à un état $\psi$). Pour chaque propriété sont
détaillés sa signification opérationnelle, sa formalisation, l'algorithme
de vérification appliqué et l'interprétation du résultat.

\subsubsection*{Algorithme générique de vérification LTL utilisé}

L'analyseur n'implémente pas un model checker LTL générique (à base
d'automates de Büchi). Il évalue les propriétés via trois primitives
adaptées à la structure du graphe d'accessibilité :

\begin{enumerate}[leftmargin=1.2cm,itemsep=2pt]
  \item \textbf{$G(\varphi)$} : on vérifie $\varphi$ dans \emph{chaque}
        marquage de $\mathrm{Reach}(M_0)$. Premier contre-exemple $\Rightarrow$
        violation.
  \item \textbf{$F(\varphi)$} (à partir d'un marquage $M$) : BFS depuis $M$
        ; on cherche un descendant satisfaisant $\varphi$.
  \item \textbf{$G(\varphi \rightarrow F(\psi))$} : pour chaque marquage
        satisfaisant $\varphi$, on lance le test $F(\psi)$ ci-dessus.
        Échec sur l'un d'eux $\Rightarrow$ violation.
\end{enumerate}

L'absence de cycles infinis sans progression est assurée par la finitude
de l'espace d'états (111 marquages).

% ----------- LTL1 -----------
\subsubsection*{LTL1 -- Exclusion temporelle tir/rechargement (sûreté)}

\paragraph{Énoncé naturel.}
« À aucun instant le système ne tire et ne recharge simultanément. »

\paragraph{Formule.}
\[
G\bigl(\neg(\textit{firing} \wedge \textit{reloading})\bigr)
\]
où $\textit{firing} \equiv M(P_6) \geq 1$ et
$\textit{reloading} \equiv M(P_7) \geq 1$.

\paragraph{Vérification.}
Évaluation directe de $G$ : pour chaque marquage atteignable, on teste
$\neg(M(P_6) \geq 1 \wedge M(P_7) \geq 1)$. La preuve structurelle
provient également du P-invariant $y_1$ (cf.\ INV5).

\paragraph{Garantie apportée.}
\OK\ -- vérifiée. C'est la version \emph{temporelle} de l'invariant INV5 :
là où INV5 affirme la cohérence ponctuelle de chaque état, LTL1 la
\emph{quantifie sur le temps}, ce qui correspond à l'invariant tel qu'un
opérateur le percevrait dans un journal de bord.

% ----------- LTL2 -----------
\subsubsection*{LTL2 -- Tout tir conduit à une journalisation (vivacité)}

\paragraph{Énoncé naturel.}
« Si le système tire, alors il existe nécessairement un état futur où
l'événement est journalisé. »

\paragraph{Formule.}
\[
G\bigl(\textit{fire} \;\Longrightarrow\; F(\textit{log\_recorded})\bigr)
\]
avec $\textit{fire} \equiv M(P_6) \geq 1$ et
$\textit{log\_recorded} \equiv M(P_{12}) \geq 1$.

\paragraph{Vérification.}
Pour chaque marquage avec $M(P_6) \geq 1$, BFS depuis ce marquage à la
recherche d'un descendant avec $M(P_{12}) \geq 1$. La transition $T_5$
produit systématiquement un jeton sur $P_{11}$, et $T_{11}$ est
\emph{toujours} franchissable lorsque $P_{11} \geq 1$ (aucune contention
sur cette transition).

\paragraph{Garantie apportée.}
\OK\ -- vérifiée. La traçabilité du tir est garantie sur tous les chemins
d'exécution, condition sine qua non d'un audit Kafka complet. C'est l'analogue
LTL de l'invariant INV8.

% ----------- LTL3 -----------
\subsubsection*{LTL3 -- Récupération automatique après erreur}

\paragraph{Énoncé naturel.}
« Toute apparition d'erreur dans le système est suivie par un retour à
l'état de repos \textit{Idle} en un nombre fini d'étapes. »

\paragraph{Formule.}
\[
G\bigl(\textit{error} \;\Longrightarrow\; F(\textit{idle})\bigr)
\]
avec $\textit{error} \equiv M(P_9) \geq 1$ et
$\textit{idle} \equiv M(P_0) \geq 1$.

\paragraph{Vérification.}
Pour chaque marquage avec $M(P_9) \geq 1$, on cherche un descendant
satisfaisant $M(P_0) \geq 1$. La transition $T_{10}$ ($P_9 \to P_0$) est
toujours franchissable depuis tout marquage où $P_9$ est marquée. Le
T-invariant $\{T_9, T_{10}\}$ confirme que le cycle erreur/récupération est
réversible.

\paragraph{Garantie apportée.}
\OK\ -- vérifiée. Cette propriété est essentielle pour la \emph{tolérance
aux pannes} : le système ne peut s'enliser indéfiniment dans un état
d'erreur, même en cas de fautes répétées. C'est l'analogue Petri Net de la
stratégie \code{Restart} du \code{SupervisorActor} Akka.

% ----------- LTL4 -----------
\subsubsection*{LTL4 -- Vivacité globale (\emph{global liveness})}

\paragraph{Énoncé naturel.}
« Quel que soit l'état courant, l'état \textit{Idle} reste toujours
atteignable dans le futur. »

\paragraph{Formule.}
\[
G\bigl(F(\textit{idle})\bigr)
\]
soit, plus explicitement, $\forall M \in \mathrm{Reach}(M_0),\
\exists\, M' \in \mathrm{Reach}(M),\ M'(P_0) \geq 1$.

\paragraph{Vérification.}
Pour chaque marquage, BFS sortant à la recherche de $M(P_0) \geq 1$. Cette
propriété est plus forte que LTL3 (récupération conditionnelle) : elle
exige le retour à \textit{Idle} \emph{depuis tout état}, pas seulement
depuis les états d'erreur.

\paragraph{Garantie apportée.}
\OK\ -- vérifiée. C'est la formalisation LTL de la propriété de
\emph{réinitialisation} INV9. Elle garantit que le système ne possède pas
d'état « absorbant » dont on ne peut sortir, à l'exception sémantiquement
attendue de l'épuisement complet du stock (où le système atteint encore
\textit{Idle} mais sans pouvoir tirer à nouveau, ce qui est correct).

% ----------- LTL5 -----------
\subsubsection*{LTL5 -- Conservation positive du stock (sûreté)}

\paragraph{Énoncé naturel.}
« Le stock de munitions reste à tout instant positif ou nul ; aucun
\emph{underflow} n'est possible. »

\paragraph{Formule.}
\[
G\bigl(\neg(\textit{ammo} < 0)\bigr) \;\equiv\; G\bigl(M(P_{10}) \geq 0\bigr).
\]

\paragraph{Vérification.}
Conséquence immédiate de la sémantique des réseaux de Petri : un jeton ne
peut être consommé que s'il existe. La propriété est confirmée par
énumération sur les 111 marquages.

\paragraph{Garantie apportée.}
\OK\ -- vérifiée. C'est la version LTL de INV4. L'intérêt d'exprimer la
même contrainte en logique temporelle est de permettre des combinaisons
ultérieures, par exemple
$G(\textit{ammo} = 0 \rightarrow \neg F(\textit{firing}))$
(« plus de munition implique plus jamais de tir »), qui est elle aussi
implicitement satisfaite.

% ----------- LTL6 -----------
\subsubsection*{LTL6 -- Exclusion temporelle cooldown/firing}

\paragraph{Énoncé naturel.}
« Pendant la phase de \textit{cooldown}, aucun tir ne peut se produire. »

\paragraph{Formule.}
\[
G\bigl(\textit{cooldown} \;\Longrightarrow\; \neg\,\textit{firing}\bigr)
\]
avec $\textit{cooldown} \equiv M(P_8) \geq 1$.

\paragraph{Vérification.}
Évaluation directe : pour chaque marquage atteignable, test de
$M(P_8) \geq 1 \Rightarrow M(P_6) = 0$. Garantie structurelle par le
P-invariant $y_1$ (les places de contrôle se partagent un unique jeton).

\paragraph{Garantie apportée.}
\OK\ -- vérifiée. C'est la formalisation LTL de INV6. L'analyseur ne
trouve aucun marquage où $P_6$ et $P_8$ seraient simultanément marquées.

% ----------- Synthèse LTL -----------
\subsubsection*{Synthèse des propriétés LTL}

\begin{table}[H]
\centering\small
\begin{tabularx}{\linewidth}{c X l c}
\toprule
\rowcolor{tableheader}\textbf{\#} & \textbf{Formule} & \textbf{Type} & \textbf{Résultat} \\
\midrule
LTL1 & $G(\neg(\textit{firing} \wedge \textit{reloading}))$       & Sûreté            & \OK \\
LTL2 & $G(\textit{fire} \rightarrow F(\textit{log\_recorded}))$    & Vivacité          & \OK \\
LTL3 & $G(\textit{error} \rightarrow F(\textit{idle}))$            & Récupération      & \OK \\
LTL4 & $G(F(\textit{idle}))$                                       & Vivacité globale  & \OK \\
LTL5 & $G(\neg(\textit{ammo} < 0))$                                & Sûreté            & \OK \\
LTL6 & $G(\textit{cooldown} \rightarrow \neg \textit{firing})$     & Sûreté            & \OK \\
\bottomrule
\end{tabularx}
\caption{Synthèse de la vérification des propriétés LTL.}
\end{table}

\paragraph{Discussion.}
Les propriétés LTL1, LTL5, LTL6 forment le \emph{bloc sûreté} : elles
ferment la porte à six classes de comportements indésirables (tir
concurrent au rechargement, sous-stock négatif, tir pendant le cooldown,
\dots). Les propriétés LTL2 et LTL3 forment le \emph{bloc vivacité} : elles
garantissent que les progrès attendus (journalisation, récupération
d'erreur) finissent toujours par se produire. LTL4 est une propriété de
\emph{vivacité globale} qui interdit toute forme de famine ou
d'enlisement. Conjointement avec les invariants métier, elles caractérisent
formellement le comportement nominal acceptable du \fcs.

\noindent
\textbf{Bilan : 6/6 propriétés LTL vérifiées.} Combinées aux invariants
métier, elles forment un \emph{filet de garanties formelles} couvrant à la
fois les états instantanés et les évolutions temporelles du système.

\subsection{Récapitulatif global}

\begin{table}[H]
\centering\small
\begin{tabular}{l c}
\toprule
\rowcolor{tableheader}\textbf{Catégorie} & \textbf{Résultat} \\
\midrule
Invariants métier  & 10/10 \PASS \\
Propriétés LTL     & 6/6 \PASS \\
P-invariants       & 2 trouvés, valides \PASS \\
T-invariants       & 1 trouvé, valide \PASS \\
Bornitude          & 6-borné \PASS \\
Vivacité           & $L_1$/$L_4$ (attendu) \PASS \\
Deadlocks          & 0 \PASS \\
\bottomrule
\end{tabular}
\caption{Synthèse de la vérification formelle.}
\end{table}

\paragraph{Conclusion.}
Le modèle formel satisfait toutes les propriétés de sûreté et de vivacité
attendues. Le système est \emph{exempt de deadlocks} et respecte tous les
invariants métier identifiés.

% ============================================================
\section{Simulation comparée : Petri Net vs Akka}
% ============================================================

\subsection{Méthodologie}

La simulation comparée vise à vérifier que le \emph{modèle formel} (réseau
de Petri) et l'\emph{implémentation} (Akka) produisent des comportements
cohérents.

\begin{enumerate}[leftmargin=1.2cm,itemsep=2pt]
  \item Exécution déterministe de transitions sur le Petri Net avec
        enregistrement de la trace ;
  \item Exécution du système Akka avec les mêmes scénarios ;
  \item Comparaison via le mapping \emph{transition $\leftrightarrow$
        événement Akka}.
\end{enumerate}

\subsubsection{Mapping Transition $\to$ Événement Akka}

\begin{table}[H]
\centering\small
\begin{tabularx}{\linewidth}{l X}
\toprule
\rowcolor{tableheader}\textbf{Transition} & \textbf{Événement Akka} \\
\midrule
$T_0$ \code{detect\_target}     & \code{SensorActor} $\to$ \code{TargetDetected} $\to$ \code{InitiateFireCycle} \\
$T_1$ \code{lock\_target}       & \code{TrackingActor} $\to$ \code{TargetLockConfirmed} \\
$T_2$ \code{load\_ammo}         & \code{AmmoActor} $\to$ \code{AmmoLoadConfirmed} \\
$T_3$ \code{authorize\_fire}    & \code{CommandActor} $\to$ \code{FireAuthConfirmed} \\
$T_4$ \code{ready\_sync}        & \code{FireControlActor.checkSync} $\to$ \code{readyToFire} \\
$T_5$ \code{fire}               & \code{FireControlActor.executeFire} $\to$ \code{FireExecuted} \\
$T_6$ \code{end\_fire}          & \code{FireControlActor} $\to$ \code{postFire} \\
$T_7$ \code{reload\_complete}   & \code{FireControlActor} $\to$ \code{ReloadTimeout} \\
$T_8$ \code{cooldown\_complete} & \code{FireControlActor} $\to$ \code{CooldownTimeout} $\to$ \code{Idle} \\
$T_9$ \code{error\_detected}    & \code{SupervisorActor} $\to$ \code{ReportError} \\
$T_{10}$ \code{error\_recovery} & \code{SupervisorStrategy.restart} $\to$ \code{Idle} \\
$T_{11}$ \code{kafka\_log}      & \code{KafkaProducerActor} $\to$ \code{PublishEvent} \\
\bottomrule
\end{tabularx}
\caption{Mapping entre transitions du réseau de Petri et événements Akka.}
\end{table}

\subsection{Résultats par scénario}

\subsubsection{Scénario 1 : cycle nominal}

Séquence : $T_0 \to T_1 \to T_2 \to T_3 \to T_4 \to T_5 \to T_6 \to T_7 \to
T_8 \to T_{11} \to T_{11}$.

\noindent
Toutes les transitions s'exécutent sans blocage. Le système retourne à
\textit{Idle} avec deux munitions restantes et deux événements journalisés.
Le cycle Akka suit la même séquence.\\
\textbf{Verdict : \PASS\ -- CONFORME.}

\subsubsection{Scénario 2 : tir sans autorisation}

$T_0 \to T_1 \to T_2 \to T_4$ (avec $T_3$ omis). $T_4$ \emph{bloque} car
$P_4$ (\code{Fire\_Authorized}) est vide. En Akka, \code{CommandActor}
renvoie \code{FireAuthDenied} ; le cycle est avorté.\\
\textbf{Verdict : \PASS\ -- CONFORME.}

\subsubsection{Scénario 3 : tir sans munition chargée}

$T_0 \to T_1 \to T_3 \to T_4$ (avec $T_2$ omis). $T_4$ \emph{bloque} car
$P_3$ (\code{Ammo\_Loaded}) est vide. En Akka, \code{AmmoActor} renvoie
\code{AmmoLoadFailed}.\\
\textbf{Verdict : \PASS\ -- CONFORME.}

\subsubsection{Scénario 4 : erreur et récupération}

$T_9 \to T_{10}$. Le cycle s'exécute, retour à \textit{Idle}. En Akka,
stratégie \emph{restart} du \code{SupervisorActor}.\\
\textbf{Verdict : \PASS\ -- CONFORME.}

\subsubsection{Scénario 5 : épuisement des munitions}

Cycle complet avec un stock initial $\textit{ammo} = 1$, puis tentative de
nouveau cycle. $T_0$ \emph{bloque} car $P_{10}=0$. En Akka,
\code{AmmoActor} renvoie \code{AmmoLoadFailed}.\\
\textbf{Verdict : \PASS\ -- CONFORME.}

\subsubsection{Scénario 6 : double tir consécutif}

Après $T_5$, tentative immédiate d'un nouveau $T_5$. La transition
\emph{bloque} car $P_5$ est vide. En Akka, \code{FireControlActor} se
trouve en état \emph{postFire} et n'accepte plus l'événement.\\
\textbf{Verdict : \PASS\ -- CONFORME.}

\subsection{Récapitulatif}

\begin{table}[H]
\centering\small
\begin{tabular}{l c c c}
\toprule
\rowcolor{tableheader}\textbf{Scénario} & \textbf{Petri Net} & \textbf{Akka} & \textbf{Verdict} \\
\midrule
Cycle nominal           & Succès        & Succès        & CONFORME \\
Sans autorisation       & Blocage $T_4$ & Abort         & CONFORME \\
Sans munition           & Blocage $T_4$ & Abort         & CONFORME \\
Erreur/récupération     & Succès        & Succès        & CONFORME \\
Épuisement munitions    & Blocage $T_0$ & Abort         & CONFORME \\
Double tir              & Blocage $T_5$ & Impossible    & CONFORME \\
\bottomrule
\end{tabular}
\caption{Résultats des six scénarios de simulation comparée.}
\end{table}

\noindent
\textbf{Résultat : 6/6 scénarios conformes.}

\subsection{Analyse}

\paragraph{Différences structurelles.}
\begin{enumerate}[leftmargin=1.2cm,itemsep=2pt]
  \item \textbf{Concurrence} : en Akka, $T_1$, $T_2$ et $T_3$ sont
        véritablement concurrents (messages asynchrones). Dans le réseau
        de Petri, ils sont entrelacés mais le résultat est identique grâce
        à la synchronisation $T_4$.
  \item \textbf{Temporisation} : Akka utilise des \emph{timers} réels
        (2\,s pour le \emph{reload}, 3\,s pour le \emph{cooldown}). Le
        réseau de Petri est non temporisé.
  \item \textbf{Granularité} : $T_0$ vérifie $P_{10} \geq 1$ alors qu'en
        Akka c'est l'\code{AmmoActor} qui vérifie le stock lors du
        \code{LoadAmmo}. Le modèle formel est plus conservatif.
\end{enumerate}

\paragraph{Conclusion.}
Le modèle formel constitue une \emph{abstraction fidèle} de l'implémentation.
Les propriétés de sûreté prouvées sur le réseau de Petri s'appliquent au
système Akka, sous réserve que l'implémentation respecte les contrats
modélisés.

\subsection{Commandes SBT du projet}

Le projet est piloté par \textbf{sbt} (Scala Build Tool).

\paragraph{\code{sbt compile}.}
Compile les sources Scala (\code{src/main/scala/}). Cette commande :
\begin{itemize}[leftmargin=1.2cm,itemsep=1pt]
  \item résout et télécharge les dépendances déclarées dans \code{build.sbt}
        (Akka Typed, Kafka, ScalaTest, etc.) ;
  \item compile les 8 acteurs Akka (\code{actors/}), les modèles de
        messages (\code{model/}), le réseau de Petri (\code{petri/}) et le
        point d'entrée (\code{Main.scala}) ;
  \item produit les fichiers \code{.class} dans \code{target/}.
\end{itemize}

\paragraph{\code{sbt test}.}
Exécute l'ensemble des tests unitaires et de vérification formelle :

\begin{table}[H]
\centering\small
\begin{tabularx}{\linewidth}{l l X}
\toprule
\rowcolor{tableheader}\textbf{Suite} & \textbf{Fichier} & \textbf{Vérifie} \\
\midrule
\code{ActorSpec}            & \code{actors/ActorSpec.scala}
                            & Comportement des 8 acteurs (envoi/réception, transitions, supervision). \\
\code{PetriNetSpec}         & \code{petri/PetriNetSpec.scala}
                            & Moteur générique du réseau (franchissement, pré/post, marquage). \\
\code{FCSPetriNetSpec}      & \code{petri/FCSPetriNetSpec.scala}
                            & Réseau spécifique au \fcs (13 places, 12 transitions, scénarios). \\
\code{VerificationSpec}     & \code{petri/VerificationSpec.scala}
                            & 10 invariants métier, 6 propriétés LTL, exploration, deadlocks. \\
\code{InvariantAnalysisSpec}& \code{petri/InvariantAnalysisSpec.scala}
                            & Analyse structurelle (P/T-invariants, bornitude, vivacité). \\
\bottomrule
\end{tabularx}
\caption{Suites de tests du projet.}
\end{table}

Cette commande ne nécessite \emph{pas} de serveur Kafka.

\paragraph{\code{sbt "run akka-demo"}.}
Lance une démonstration interactive du système Akka/Kafka :
\begin{enumerate}[leftmargin=1.2cm,itemsep=1pt]
  \item création d'un \code{ActorSystem} avec le \code{SupervisorActor} en
        racine ;
  \item démarrage (\code{StartSystem}) : les 8 acteurs sont créés et
        supervisés ;
  \item exécution d'un cycle nominal : détection $\to$ verrouillage $\to$
        chargement $\to$ autorisation $\to$ synchronisation $\to$ tir $\to$
        rechargement $\to$ refroidissement ;
  \item journalisation Kafka (mode local en l'absence de broker) ;
  \item arrêt propre par appui sur \emph{Entrée}.
\end{enumerate}

\paragraph{\code{sbt "run conformance"}.}
Exécute la \emph{vérification de conformité} entre modèle formel et
implémentation, en trois phases :
\begin{enumerate}[leftmargin=1.2cm,itemsep=1pt]
  \item \textbf{Simulation Petri Net} : 6 scénarios sur le réseau,
        enregistrement des traces ;
  \item \textbf{Simulation Akka} : mêmes scénarios dans des
        \code{ActorSystem} dédiés, collecte des événements
        \code{KafkaProducerActor} ;
  \item \textbf{Comparaison} : pour chaque scénario, application du
        mapping et production d'un rapport CONFORME / NON-CONFORME.
\end{enumerate}

% ============================================================
\section{Mode live interactif}
% ============================================================

\subsection{Objectif}

Le mode \code{live} ajoute un \textbf{panneau de commande interactif} pour
piloter le système pas à pas, observer l'état courant et distinguer
clairement les actions \textbf{manuelles} (opérateur) des actions
\textbf{automatiques} (système). Il complète les exécutions
\code{akka-demo} et \code{conformance} avec une vue orientée
exploitation/formation.

\subsection{Architecture technique}

Le mode live repose sur deux points d'entrée :

\begin{enumerate}[leftmargin=1.2cm,itemsep=2pt]
  \item \code{Main.runInteractive(modeArg: Option[String])} :
  \begin{itemize}[leftmargin=1.0cm,itemsep=1pt]
    \item construit le réseau via \code{FCSPetriNet.build} ;
    \item choisit le mode d'affichage (\emph{verbose} par défaut,
          \emph{compact} optionnel) ;
    \item délègue l'exécution à \code{InteractiveSimulator.run}.
  \end{itemize}
  \item \code{InteractiveSimulator} :
  \begin{itemize}[leftmargin=1.0cm,itemsep=1pt]
    \item gère la boucle interactive ;
    \item classe les transitions (manuelles vs automatiques) ;
    \item applique les délais simulés ;
    \item affiche l'état système et les actions disponibles.
  \end{itemize}
\end{enumerate}

\subsection{Fonctionnalités ajoutées}

\paragraph{Modes d'affichage.}
\begin{itemize}[leftmargin=1.2cm,itemsep=1pt]
  \item \textbf{Verbose} (par défaut) : journal détaillé, transitions
        automatiques explicitées, compte à rebours avant tir.
  \item \textbf{Compact} : sortie concise, invite simplifiée, message
        \code{Waiting for operator action...}.
\end{itemize}

\paragraph{Ergonomie de conduite.}
\begin{itemize}[leftmargin=1.2cm,itemsep=1pt]
  \item Affichage structuré en sections \code{STATUS}, \code{AUTO ACTIONS},
        \code{MANUAL ACTIONS}, \code{INPUT}.
  \item Commandes opérateur : numéro d'action, \code{status}, \code{help},
        \code{reset}, \code{quit}.
\end{itemize}

\paragraph{Comportement fonctionnel.}
\begin{itemize}[leftmargin=1.2cm,itemsep=1pt]
  \item Progression automatique des transitions non manuelles.
  \item Délais simulés (chargement, autorisation, cooldown, etc.).
  \item Compte à rebours de tir : \code{3}, \code{2}, \code{1},
        \code{Fire}.
  \item Avertissement de stock faible dédupliqué (pas de \emph{spam}).
  \item Comptage de cycles incrémenté uniquement au retour valide à
        \emph{Idle}.
\end{itemize}

\subsection{Commandes SBT associées}

\begin{table}[H]
\centering\small
\begin{tabular}{l l}
\toprule
\rowcolor{tableheader}\textbf{Commande} & \textbf{Effet} \\
\midrule
\code{sbt "run live"}            & Panneau interactif en mode \emph{verbose} (défaut). \\
\code{sbt "run live compact"}    & Panneau interactif en mode \emph{compact}. \\
\code{sbt "run live verbose"}    & Mode \emph{verbose} explicite. \\
\bottomrule
\end{tabular}
\caption{Commandes pour lancer le mode live.}
\end{table}

\subsection{Validation et tests}

La fonctionnalité \emph{live} est couverte par \code{InteractiveSimulatorSpec} :
classification des transitions, lisibilité des labels, parsing du mode
(\emph{verbose}/\emph{compact}/inconnu), présence conditionnelle du message
d'attente, séquence du compte à rebours, délais, progression automatique
sur scénarios nominal et erreur, cas limites (munitions nulles, cycles
multiples). Le résultat attendu en CI/local est une suite verte couvrant
l'ensemble des tests du projet.

% ============================================================
\section{Conclusion}
% ============================================================

Ce projet a permis de mener à bien la chaîne complète de modélisation,
implémentation et vérification d'un système distribué critique :

\begin{itemize}[leftmargin=1.2cm,itemsep=2pt]
  \item une \textbf{architecture Akka/Scala} en 8 acteurs typés couvre la
        détection, le suivi balistique, la gestion des munitions,
        l'autorisation et l'orchestration du tir ;
  \item un \textbf{réseau de Petri} de 13 places et 12 transitions
        capture les états et transitions critiques ;
  \item un \textbf{analyseur entièrement maison} (sans outil tiers) calcule
        les invariants par élimination de Gauss et explore l'espace
        d'états par BFS ;
  \item l'\textbf{exploration exhaustive} (111 marquages, 227 arcs)
        confirme l'\emph{absence totale de deadlock} ;
  \item \textbf{10 invariants métier} et \textbf{6 propriétés LTL} de
        sûreté/vivacité sont vérifiés ;
  \item six \textbf{scénarios de simulation comparée} confirment la
        cohérence entre le modèle formel et l'implémentation Akka ;
  \item un \textbf{mode interactif} permet le pilotage manuel et
        l'observation pas à pas du système.
\end{itemize}

L'apport pédagogique principal réside dans la combinaison de deux
approches complémentaires : \emph{model checking} pour la rigueur
mathématique (preuves d'absence de deadlock et de respect des invariants)
et \emph{simulation Akka} pour le réalisme opérationnel.

\paragraph{Liens du projet.}
Le code source, les tests et la documentation complète sont disponibles à
l'adresse :
\begin{center}
\url{https://github.com/mcrayssac/Fire-Control-System}
\end{center}

% ============================================================
\section*{Bibliographie}
\addcontentsline{toc}{section}{Bibliographie}
% ============================================================

\begin{thebibliography}{99}

\bibitem{petri1962}
Petri, C.~A. (1962).
\emph{Kommunikation mit Automaten}.
Thèse de doctorat, Université de Bonn.

\bibitem{murata1989}
Murata, T. (1989).
\emph{Petri Nets: Properties, Analysis and Applications}.
Proceedings of the IEEE, 77(4), 541--580.

\bibitem{baier2008}
Baier, C. \& Katoen, J.-P. (2008).
\emph{Principles of Model Checking}.
MIT Press.

\bibitem{clarke1999}
Clarke, E.~M., Grumberg, O. \& Peled, D.~A. (1999).
\emph{Model Checking}.
MIT Press.

\bibitem{pnueli1977}
Pnueli, A. (1977).
\emph{The Temporal Logic of Programs}.
18th Annual Symposium on Foundations of Computer Science (FOCS), IEEE.

\bibitem{hewitt1973}
Hewitt, C., Bishop, P. \& Steiger, R. (1973).
\emph{A Universal Modular ACTOR Formalism for Artificial Intelligence}.
IJCAI'73.

\bibitem{agha1986}
Agha, G. (1986).
\emph{Actors: A Model of Concurrent Computation in Distributed Systems}.
MIT Press.

\bibitem{lightbend2024}
Lightbend (2024).
\emph{Akka Documentation}.
\url{https://doc.akka.io/}

\bibitem{lions1996ariane}
Lions, J.-L. \emph{et al.} (1996).
\emph{Ariane 5 Flight 501 Failure Report}.
ESA/CNES Inquiry Board.

\bibitem{peterson1981}
Peterson, J.~L. (1981).
\emph{Petri Net Theory and the Modeling of Systems}.
Prentice-Hall.

\end{thebibliography}

\end{document}
